% !TEX root = main2.tex

\section{Invariant verification}
\label{sec:reachalgo}

A subproblem for invariant verification is to
compute $\reachtube{\H}$, or more specifically, the reachtubes for the set of trajectories $\TL$ in a given mode, up to a time bound.
%
This is a difficult problem, even when $\TL$ is generated by white-box models.
%
The algorithms in~\cite{donze2010breach,DMV:EMSOFT2013,FanMitra:2015df} approximate reachtubes using simulations and sensitivity analysis of ODE models generating $\TL$.
%
Here, we begin with a probabilistic method for estimating sensitivity from black-box simulators.
 

\subsection{Discrepancy functions}
\label{sec:disc}

Sensitivity of trajectories is formalized by the notion of discrepancy
functions \cite{DMV:EMSOFT2013}.  For a set 
$\TL$, a {\em discrepancy function\/} is a uniformly continuous
function $\beta: \reals^n \times \reals^n \times \nnreals \rightarrow
\nnreals$, such that for any pair of identically labeled trajectories
$\langle \tau_1,\ell \rangle, \langle \tau_2, \ell \rangle \in \TL$,
and any $t \in \tau_1.\dom \cap \tau_2.\dom$:
\begin{inparaenum}[(a)]
	\item $\beta$ upper-bounds the distance between the trajectories, i.e., 
	\begin{align}
	|\tau_1(t) - \tau_2(t)|  \leq \beta(\tau_1.\fstate,\tau_2.\fstate,t), \label{eq:discrepancy}
	\end{align} and
	\item $\beta$ converges to $0$ as the initial states converge, i.e., for any trajectory $\tau$ and $t \in
          \tau.\dom$, if a sequence of trajectories $\tau_1,\ldots,
          \tau_k, \ldots$ has $\tau_k.\fstate \rightarrow
          \tau.\fstate$, then $\beta(\tau_k.\fstate,$
          $\tau.\fstate,t)$ $\rightarrow 0$.
\end{inparaenum}
In~\cite{DMV:EMSOFT2013} it is shown how given a $\beta$, condition~(a) can used to over-approximate reachtubes  from simulations, 
and condition~(b) can be used to make these approximations arbitrarily precise. 
%
Techniques for computing $\beta$ from ODE models are developed in~\cite{FanMitra:2015df,FanM:EMSOFT2016,HFMMK:CAV2014}, but these are not applicable here in absence of such models. Instead we present a simple method for discovering discrepancy functions that only uses simulations. Our method is based on classical results on PAC learning linear separators~\cite{KearnsVazirani}. We recall these before applying them to find discrepancy functions.
\vspace{-10pt}
\subsubsection{Learning linear separators.}
\label{Sec: discrepancy}
For $\Gamma \subseteq \reals\times\reals$, a \emph{linear separator}
is a pair $(a,b) \in \reals^2$ such that
\begin{align}
\forall (x,y) \in \Gamma.\ x \leq ay + b.
\label{eq:separator}
\end{align}
Let us fix a subset $\Gamma$ that has a (unknown) linear separator
$(a_*,b_*)$. Our goal is to discover some $(a,b)$ that is a linear
seprator for $\Gamma$ by sampling points in $\Gamma$~\footnote{We
  prefer to present the learning question in this form as opposed to
  one where we learn a Boolean concept because it is closer to the
  task at hand.}. The assumption is that elements of $\Gamma$ can be
drawn according to some (unknown) distribution ${\cal D}$. With
respect to ${\cal D}$, the \emph{error} of a pair $(a,b)$ from
satisfying Equation~\ref{eq:separator}, is defined to be
$\mathsf{err}_{{\cal D}}(a,b) = {\cal D}(\{(x,y) \in \Gamma\: |\: x >
ay+b\})$
where ${\cal D}(X)$ is the measure of set $X$ under distribution
${\cal D}$. Thus, the error is the measure of points (w.r.t. ${\cal
  D}$) that $(a,b)$ is not a linear separator for. There is a very
simple (probabilistic) algorithm that finds a pair $(a,b)$ that is a
linear separator for a large fraction of points in $\Gamma$, as
follows.
\begin{enumerate}
\itemsep0em 
\item\label{alg:linsep1} Draw $k$ pairs $(x_1,y_1), \ldots (x_k,y_k)$
  from $\Gamma$ according to ${\cal D}$; the value of $k$ will be
  fixed later.
\item\label{alg:linsep2} Find $(a,b) \in \reals^2$ such that $x_i
  \leq ay_i + b$ for all $i \in \{1,\ldots k\}$.
\end{enumerate}
Step~\ref{alg:linsep2} involves checking feasibility of a linear
program, and so can be done efficiently. This algorithm, with high
probability, finds a linear separator for a large fraction of points.
%
\begin{proposition}
\proplabel{linear-sep-learn} Let $\epsilon, \delta \in \plreals$. If
$k \geq \frac{1}{\epsilon}\ln\frac{1}{\delta}$ then, with probability
$\geq 1-\delta$, the above algorithm finds $(a,b)$ such that
$\mathsf{err}_{{\cal D}}(a,b) < \epsilon$.
\end{proposition}
%
%The proof follows from the PAC-learnability of concepts with low
%VC-dimension \cite{KearnsVazirani}. The detailed proof is provided in the Appendix \ref{sec:proof}.
\begin{proof}
The result follows from the PAC-learnability of concepts with low
VC-dimension~\cite{KearnsVazirani}. However, since the proof is very
simple in this case, we reproduce it here for completeness. Let $k$ be
as in the statement of the proposition, and suppose the pair $(a,b)$
identified by the algorithm has error $> \epsilon$. We will bound the
probability of this happening.

Let $B = \{(x,y)\: |\: x > ay+b\}$. We know that ${\cal D}(B) >
\epsilon$. The algorithm chose $(a,b)$ only because no element from
$B$ was sampled in Step~\ref{alg:linsep1}. The probability that this
happens is $\leq (1-\epsilon)^k$. Observing that $(1-s) \leq e^{-s}$
for any $s$, we get $(1-\epsilon)^k \leq e^{-\epsilon k} \leq e^{-\ln
  \frac{1}{\delta}} = \delta$.  This gives us the desired result.
\end{proof}


\subsubsection{Learning discrepancy functions}
Discrepancy functions will be computed from simulation data
independently for each mode. Let us fix a mode $\ell \in \L$, and a
domain $[0,T]$ for each trajectory. The discrepancy functions that we
will learn from simulation data, will be one of two different forms,
and we discuss how these are obtained.

\vspace{-10pt}
\paragraph{Global exponential discrepancy (GED)} is a function of the form
\[
\beta(x_1,x_2,t) = |x_1 - x_2| Ke^{\gamma t}.
\]
Here $K$ and $\gamma$ are constants. Thus, for any pair of
trajectories $\tau_1$ and $\tau_2$ (for mode $\ell$), we have
\[
\forall t \in [0,T].\ |\tau_1(t) - \tau_2(t)| \leq |\tau_1.\fstate -
\tau_2.\fstate| Ke^{\gamma t}.
\]
Taking logs on both sides and rearranging terms, we have
\[
\forall t.\ \ln \frac{|\tau_1(t) - \tau_2(t)|}{|\tau_1.\fstate -
\tau_2.\fstate|} \leq \gamma t + \ln K.
\]
It is easy to see that a global exponential discrepancy is nothing but
a linear separator for the set $\Gamma$ consisting of pairs $(\ln
\frac{|\tau_1(t) = \tau_2(t)|}{|\tau_1.\fstate - \tau_2.\fstate|}, t)$
for all pairs of trajectories $\tau_1,\tau_2$ and time $t$. Using the
sampling based algorithm described before, we could construct a GED for a mode $\ell \in \L$, where sampling from
$\Gamma$ reduces to using the simulator to generate traces from
different states in $\TL_{\sf init, \ell}$. \propref{linear-sep-learn}
guarantees the correctness, with high probability, for any separator
discovered by the algorithm.  However, for our reachability algorithm
to not be too conservative, we need $K$ and $\gamma$ to be
small. Thus, when solving the linear program in Step~\ref{alg:linsep2}
of the algorithm, we search for a solution minimizing  $\gamma T +
\ln K$.

\vspace{-10pt}
\paragraph{Piece-wise exponential discrepancy (PED).} 
The second form of discrepancy functions we consider, depends upon
dividing up the time domain $[0,T]$ into smaller intervals, and
finding a global exponential discrepancy for each interval. Let $0 =
t_0,t_1,\ldots t_N = T$ be an increasing sequence of time points. Let
$K, \gamma_1, \gamma_2, \ldots \gamma_N$ be such that for every pair
of trajectories $\tau_1,\tau_2$ (of mode $\ell$), for every $i \in
\{1,\ldots, N\}$, and $t \in [t_{i-1},t_i]$, $|\tau_1(t) = \tau_2(t)|
\leq |\tau_1(t_{i-1}) - \tau_2(t_{i-1})| Ke^{\gamma_i t}$.  Under such
circumstances, the discrepancy function itself can be seen to be given
as
\[
\beta(x_1,x_2,t) = |x_1 - x_2| Ke^{\sum_{j=1}^{i-1}\gamma_j(t_j -
  t_{j-1}) + \gamma_i (t-t_{i-1})} \qquad \mbox{for } t \in [t_{i-1},t_i].
\]
If the time points $0 = t_0,t_1,\ldots t_N = T$ are fixed, then the
constants $K, \gamma_1, \gamma_2, \ldots $ $\gamma_N$ can be discovered
using the learning approach described for GED; here, to discover $\gamma_i$, we take $\Gamma_i$ to be the pairs obtained by restricting the trajectories to be between times
$t_{i-1}$ and $t_i$. The sequence of time points $t_i$ are also
dynamically constructed by our algorithm based on the following
approach. Our experience suggests that a value for $\gamma$ that is
$\geq 2$ results in very conservative reach tube
computation. Therefore, the time points $t_i$ are constructed
inductively to be as large as possible, while ensuring that $\gamma_i <
2$.

%\mitras{Not sure if we need the following psedudocode.}
%\paragraph{Chuchu's text}
%%\chuchu{This part in red is from one discussion with Mahesh: Suppose for a system the actual discrepancy function is D, and we construct a discrepancy function D' after looking at m samples. Since we pick D' to be sound on all the sample points, D' serves as a lower bound on D. Let us define Bad to be the set of points for which D' is not the discrepancy function; so the points in Bad diverge more than D' claims.}
%
%\chuchu{This could happen only because the sample never contained points from Bad. Suppose Bad is of size $> e$. The probability that we never get a point in Bad is at most $(1 - e)^m$. We want to make sure that this chance is small; say smaller than $\delta$. Then we have $(1 - e)^m < \delta$, which happens if we take $m = O(1/e(log \delta))$ samples.
%So if we take $O(1/e(log \delta))$ samples and construct the discrepancy function then the chances that the constructed discrepancy function does not work for more e fraction of the points is at most $\delta$.}
%
%
%Given $m$ simulations at discrete time points $t_0,t_1,\dots,t_N=T$: $\xi(x_i,t_j)$, $i = 1,\dots,m$ and $j = 0,\dots,N$, we are using two different kinds of discrepancy functions
%\begin{enumerate}
%\item Global discrepancy $| \xi(x,t) - \xi(x',t) | \leq |x-x'| K e^{\gamma t}$ and, 
%\item Piecewise exponential discrepancy $$| \xi(x,t) - \xi(x',t) | \leq |x-x'| K e^{\sum_{i = 1}^{j-1} \gamma_i(t_i-t_{i-1}) + \gamma_j(t-t_j)}, t \in [t_j,t_{j+1}].$$
%\end{enumerate}
%
%For the global discrepancy function, the algorithm is transferred to solve the following optimization problem: 
%
%\[ \min \int_{0}^{T} K e^{\gamma t} dt \]
%\[ s.t, \forall x,x' \in \{x_1,\dots,x_m\}, t_j  \in \{t_0,\dots,t_N\}, \]
%\[| \xi(x,t_j) - \xi(x',t_j) | \leq |x-x'| K e^{\gamma (t_j - t_0)}, \]
%
%which can be furthur simplified to solve the following linear programming problem:
% 
%\[ \min N \ln(K) + \gamma \sum_{i=1}^{t}t_i \]
%\[ s.t, \forall x,x' \in \{x_1,\dots,x_m\}, t_j  \in \{t_0,\dots,t_N\}, \]
%\[ \ln (| \xi(x,t_j) - \xi(x',t_j) |)- \ln( |x-x'|) \leq  \ln(K) +  {\gamma (t_j - t_0)}, \]
%
% 
%For the piece-wise exponential discrepancy function, we are given a $K$, the only task is to compute $\gamma_i$s. 
%
%%In this algorithm, we require that the time steps are fixed. That is $\forall i = 1,\dots,N, t_i - t_{i-1} = \Delta t$
%
%The algorithm is separated into two steps:
%
%Step 1: compute the time intervals, to each of which a $\gamma_i$ is assigned. Basically, we want that in each time interval, the trajectories are not diverging too much.
%\begin{algorithm}[h!]
%\SetKwInOut{Input}{input}
%\SetKwInOut{Initially}{initially}
%\Input {Simulations $\xi(x_i,t_j)$, $i = 1,\dots,m$ and $j = 0,\dots,N$}
%\Initially{$\TimeIntvs \gets t_0$, $t_s = t_0$ }
%\While {$t_s \leq t_N$}
%{	
%	\For {$t \in \{ t_{s+1}, \dots, t_N\}$}
%		{\If {$\forall x,x' \in {x_1,\dots,x_m}, |\xi(x,t)-\xi(x',t)| > |\xi(x,t_s)-\xi(x',t_s)| e^{2(t-t_s)}$}
%		{{$t_e = t$ \;}
%		{Break \;}}
%		}
%	{$\TimeIntvs  \gets \TimeIntvs.append(t_e)$ \;}
%	{$t_s \gets t_e$ \;}
%}
%\Return $\TimeIntvs$ \;
%\end{algorithm}
%
%Step 2: For each time interval, we are going to compute the corresponding $\gamma$s inductively:
%\begin{algorithm}[h!]
%\SetKwInOut{Input}{input}
%\SetKwInOut{Initially}{initially}
%\Input {Simulations $\xi(x_i,t_j)$, $i = 1,\dots,m$ and $j = 0,\dots,N, \TimeIntvs$}
%\Initially{$IntNum \gets$ length$(\TimeIntvs)-1$ }
%\For {$i  = 1:IntNum$}
%{	
%	{$Q = \{t \in \{t_0,\dots,t_N\}~~|~~\TimeIntvs[i-1] \leq t \leq \TimeIntvs[i] \}$ \;}
%	{min $\gamma[i]$\;}
%	{s.t $\forall x,x' \in {x_1,\dots,x_m}, \forall t \in Q,$
%$| \xi(x,t) - \xi(x',t) | \leq |x-x'| K e^{\sum_{k = 1}^{i-1} \gamma[k] (\TimeIntvs[k] -\TimeIntvs[k-1]) + \gamma[i](t-\TimeIntvs[i-1])}  $ }
%}
%\Return $\gamma$ \;
%\end{algorithm}
\vspace{-10pt}
\subsubsection{Experiments on learning discrepancy}

We used the above algorithm to learn discrepancy functions for dozens
of modes with complex, nonlinear trajectories.  Our experiments
suggest that around 10-20 simulation traces are adequate for computing
both global and piece-wise discrepancy functions. For each mode we use
a set $S_{\sf train}$ of simulation traces that start from
independently drawn random initial states in $\TL_{\sf init,\ell}$ to
learn a discrepancy function.  Each trace may have $100$-$10000$ time
points, depending on the relevant time horizon and sample times.  Then
we draw another set $S_{\sf test}$ of $1000$ simulations traces for
validating the computed discrepancy. For every pair of trace in
$S_{\sf test}$ and for every time point, we check whether the computed
discrepancy satisfies Equation~\ref{eq:discrepancy}.  We 
observe that for $|S_{\sf train}| > 10$ the computed discrepancy
function is correct for $96\%$ of the points $S_{\sf test}$ in and for
$|S_{\sf train}| > 20$ it is correct for more than $99.9\%$, across all experiments.

% Resume here tomorrow
\subsection{Verification algorithm}
\label{sec:verfication_algorithm}

In this section, we present algorithms to solve the bounded
verification problem for hybrid systems using learned exponential discrepancy functions.
%
We first introduce an algorithm $\GraphReach$ (Algorithm
\ref{alg:ComputeRT}) which takes as input a hybrid system $\H =
\langle \L, \Theta, G, \TL \rangle$ and returns a set of
reachtubes---one for each vertex of $G$---such that their
union over-approximates $\reachtube{\H}$.
 % To simplify notations, we assume that $G$ has a single initial vertex $v_{\sf init}$.
 
%\mitras{This definition and notation has to be reconciled with Mahesh's definition and notation.}
%\begin{definition}
%For a set $\TL$ of trajectories, a set $S \subseteq \reals^n$, a mode $\ell \in \L$ and a finite sequence of time points $t_0,\dots,t_k = T$, a {\em reachtube} $Reachtube{}$ is a sequence of timed-stamped sets $(R_1, l, t_1), (R_2, l ,t_2), \dots , (R_k, l ,t_k)$ satisfying that $\forall \tau \in \TL$ with $\tau.\fstate \in S$, $\forall i = 1,\dots,k, \forall t \in [t_{i-1},t_{i}], \tau(t) \in R_i$. 
%\end{definition}
%$Reachtube \downarrow_{(t,t')}$ is the projection of the reachtube $RT$ into time interval $(t,t')$.

% resume here.

% As when the algorithm is used, the simulator/blackbox for the labeled trajectories $\TL$ will be clear from context, we will omit $\TL$ in the algorithm.

$\GraphReach$ maintains two  data-structures:
 \begin{inparaenum}[(a)]
\item $RS$ accumulates pairs of the form $\langle RT, v\rangle$, where  $v \in \V$ and $RT$ is its corresponding reachtube;
%%
%%
\item $\VerInit$ accumulates pairs of the form $\langle S, v \rangle$,
  where $v \in \V$ and $S\subset \reals^n$ is the set of states  from which the reachtube in $v$ is to be computed.
   %\item $RemainT$ is the remaining time. Initially, it is set to the time horizon $T$.
%\item Each vertex has an attribute $\mathrm{Time}$, which indicate the time that the algorithm should run from the given vertex. Initially, $v_{0}.\mathrm{Time}$ is set to $T$ and the rest vertices's $\mathrm{Time}$ is set to $0$.
\end{inparaenum} 
Each
$v$ could be in multiple such pairs in $RS$ and $\VerInit$.
%
Initially, $RS = \emptyset$  and
$\VerInit = \{\langle \Theta, v_{\sf init}\rangle\}$.
%, where $v_{\sf init}$ is the initial vertex of the graph $G$.


$\mathit{LearnDiscrepancy(S_{\sf init},d,\ell)}$ computes
the discrepancy function for mode $\ell$, from initial set $S_{\sf
  init}$ and upto time $d$ using the algorithm of Section \ref{sec:disc}. $\computeRT(S_{\sf init},d,\beta)$ first
generates finite simulation traces from  $S_{\sf init}$
and then bloats the traces to compute a reachtube using the discrepancy function $\beta$. This step is similar to the algorithm for dynamical systems  given in~\cite{DMV:EMSOFT2013}.
%Refering to previous work helps us to skip some of the details

The $\GraphReach$ algorithm proceeds as follows: first, a topologically sorted array of the vertices of the DAG $G$ is computed in  $\mathit{Order}$ (\lnref{ln: init2}). 
The pointer $ptr$ iterates over the $\mathit{Order}$ and for each vertex $\mathit{curv}$ the following is computed. 
%Firstly, the algorithm gets the current vertex $CurV$ and its label $l$ according to the pointer $ptr$ (\lnsref{ln: currentV}{ln: currentL}). 
The variable $\mathit{dt}$  is set to the maximum transition time to other vertices from $\mathit{curv}$ (\lnref{ln: dwt}). 
For each possible initial set $S_{\sf init}$ corresponding to $\mathit{curv}$ in $\VerInit$, the algorithm computes a discrepancy function (\lnref{ln: disc}) and uses it to compute a reachtube from $S_{\sf init}$ up to time $\mathit{dt}$ (\lnref{ln: reachtube}). 
For each successor $\mathit{nextv}$ of $\mathit{curv}$, the restriction of the computed reachtube $RT$ to the corresponding transition time interval $\edgelab((\mathit{curv,nextv}))$ is set as  an initial set for $\mathit{nextv}$ (\lnsref{ln: forloopbegin}{ln: forloopend}).
% and update $vs.\mathrm{Time}$ to be the maximum of its original value and $ CurV.\mathrm{Time} $ minus the earliest possible transition time $t$ .
%Finally, the pointer move forward to the next position (\lnref{ln: moveptr}), and the remaining time is set $0$ if current vertex is the last one in the topological sort, and is set to the next vertex's attribute $\mathrm{Time}$ otherwise.

\begin{algorithm}[h!]
\caption{$\GraphReach(\H)$ computes bounded time reachtubes for each vertex of the transition $G$ of hybrid system $\H$.}
\label{alg:ComputeRT}
\SetKwInOut{Input}{input}
\SetKwInOut{Initially}{initially}
{$RS \gets \emptyset; \VerInit \gets \{\langle \Theta, v_{\sf init} \rangle\}; \mathit{Order} \gets \mathit{TopSort}(G)$\;} \lnlabel{ln: init2}
\For {$ptr = 0: len(Order)-1$}
{	
	{$\mathit{curv} \gets \mathrm{Order}[ptr]$ \;} \lnlabel{ln: currentV}
	{$\ell \gets \vertlab(\mathit{curv})$\;} \lnlabel{ln: currentL}

	$\mathit{dt} \gets \textrm{max} \{t' \in \nnreals \:| \:\exists vs \in \V, (\mathit{curv}, vs) \in \E, (t,t') \gets \edgelab \left( (\mathit{curv}, vs) \right) \}$\; \lnlabel{ln: dwt}

	\For {$S_{\sf init} \in \{S~|~ \langle S,\mathit{curv} \rangle \in \VerInit\}$}
	{
		{$\beta \gets \mathit{LearnDiscrepancy}(S_{\sf init},\mathit{dt},\ell)$\;} \lnlabel{ln: disc}
		{$  RT \gets \computeRT(S_{\sf init},\mathit{dt},\beta)$\;} \lnlabel{ln: reachtube}
		{$RS \gets RS \cup \langle RT, \mathit{curv} \rangle $\;}
		\For {$\mathit{nextv} \in \mathit{curv}.succ$} 
		{
			$(t,t') \gets \edgelab \left( (\mathit{curv}, nextv) \right)$\; \lnlabel{ln: forloopbegin}
			$\VerInit \gets \VerInit \cup \langle  \mathit{Restr}(RT,(t,t')), nextv\rangle$\; \lnlabel{ln: forloopend}
		}
	}
}
\Return $RS$ \;
\end{algorithm}
 
%\begin{algorithm}[h!]
%\caption{Algorithm $\GraphReach(G, \Theta, T)$ to compute bounded time Reachtube for graph $G$.}
%\label{alg:ComputeRT}
%\SetKwInOut{Input}{input}
%\SetKwInOut{Initially}{initially}
%$\R \gets \emptyset; v_{0}.\mathrm{Time} \gets T; \forall v \in \V \setminus v_{\sf init}, v.\mathrm{Time} \gets 0; RemainT \gets T$\;
%{$\VerInit \gets \langle \Theta, v_{\sf init} \rangle; \mathrm{Order} \gets TopoSort(G); ptr \gets 0$\;} \lnlabel{ln: init2}
%\While {$RemainT>0$}
%{	
%	{$CurV \gets \mathrm{Order}[ptr]$ \;} \lnlabel{ln: currentV}
%	{$l \gets \vertlab(CurV)$ \;} \lnlabel{ln: currentL}
%	{$DwT \gets RemainT$ \;}\lnlabel{ln: dtimebegin}
%	\If {$CurV.successors \neq \emptyset$}
%	{
%		$Tr \gets \textrm{max} \{t' \in \nnreals | (CurV, vs) \in \E, (t,t') = \edgelab \left( (CurtV, vs) \right) \}$ \;
%		$DwT \gets \min \{Tr, RemainT\}$\; \lnlabel{ln: dtimeend}
%	}
%	\For {$S_{\sf init} \in \{S~|~ \langle S,CurV \rangle \in \VerInit\}$}
%	{
%		{$\beta \gets Discrepancy(S_{\sf init},DwT,l)$\;} \lnlabel{ln: disc}
%		{$  RT = \computeRT(S_{\sf init},DwT,\beta)$\;} \lnlabel{ln: reachtube}
%		{$\R = \R \cup \langle RT, CurV \rangle $\;}
%		\For {$vs \in CurV.successors$} 
%		{
%			$(t,t') \gets \edgelab \left( (CurtV, vs) \right)$\; \lnlabel{ln: forloopbegin}
%			$\VerInit \gets \VerInit \cup \langle  RT \downarrow_{(t,t')}, vs\rangle$\;
%			$vs.\mathrm{Time} \gets \max(vs.\mathrm{Time}, CurV.\mathrm{Time} - t)$\; \lnlabel{ln: forloopend}
%		}
%	}
%    	{$ptr \gets ptr + 1$\;} \lnlabel{ln: moveptr}
%    	{$RemainT \gets 0$\;}
%    	\If {$ptr \neq len(Order)-1$}
%    	{$RemainT \gets \mathrm{Order}[ptr].\mathrm{Time}$}
%}
%\Return $\R$ \;
%\end{algorithm}

The invariant verification algorithm $\mathit{VerifySafety}$ decides
safety of $\H$ with respect to a given unsafe set $\U$ and uses
$\GraphReach$. The detailed pseudocode appears in
Appendix~\ref{appendix:safetyveri}.  This algorithm proceeds in a way
similar to the simulation-based verification algorithms for dynamical
and hybrid systems~\cite{DMV:EMSOFT2013,fan2016automatic}.
%
Given initial set $\Theta$ and transition graph $G$ of $\H$, this algorithm partitions $\Theta$ into several subsets, and then for each subset $S$ it checks whether the computed  over-approximate reachtube $RS$ from $S$ intersects with $\U$:
 \begin{inparaenum}[(a)]
 \item If $RS$  is disjoint, the system is safe starting from $S$; 
 \item if certain part of a reachtube $RT$ is contained in $\U$, the system is declared as  unsafe and $RT$ with the the corresponding path of the graph are returned as counter-example witnesses; 
 \item if neither of the above conditions hold, then the algorithm performs refinement to get a more precise over-approximation of $RS$.
 \end{inparaenum} 
Several refinement strategies are implemented in {\toolname} to
accomplish the last step.  Broadly, these strategies rely on splitting
the initial set $S$ into smaller sets (this gives tighter discrepancy
in the subsequent vertices) and splitting the edge labels of $G$ into
smaller intervals (this gives smaller initial sets in the vertices).
 
The above description focuses on invariant properties, but the
algorithm and our implementation in {\toolname} can verify a useful
class of temporal properties.  These are properties in which the time
constraints only refer to the time since the last mode transition.
%
For example, for the $\auto{Powertrn}$ benchmark the tool verifies
requirements like ``after $4$s in $\Lmode{normal}$ mode, the air-fuel
ratio should be contained in $[14.6,14.8]$ and after $4$s in
$\Lmode{powerup}$ it should be in $[12.4,12.6]$''.

\vspace{-10pt}
\paragraph{Correctness}
Given a correct discrepancy function for each mode, we can prove the
soundness and relative completeness of
Algorithm~\ref{alg:safetyveri}. This analysis closely follows the
proof of Theorem~19 and Theorem~21
in~\cite{duggirala2015dynamic}. Combining this with the probabilistic
correctness of the $\mathit{LearnDiscrepancy}$, we obtain the
following probabilistic soundness guarantee.

%\paragraph{Soundness of algorithm}
\begin{theorem}
	If the $\beta$'s returned by $\mathit{LearnDiscrepancy}$ 
	are always  discrepancy functions for corresponding modes, then $\mathit{VerifySafety}(\H,U)$ (Algorithm~\ref{alg:safetyveri}) is sound. That is, if it outputs ``SAFE'', then $\H$  is safe with respect to $\U$ and if it outputs ``UNSAFE'' then there exists an execution of $\H$ that enters $\U$.
\end{theorem}

